\section{Discussion}
\label{sec:discussion}

\subsection{What the protocol buys, and what it costs}

The certificate turns an approximation into a statement with a hypothesis and
a conclusion: \emph{if} the interval contains exactly one eigenvalue (checked),
\emph{then} the eigenvalue lies in $[\eta-\rho^2/(b-\eta),\,\eta+\rho^2/(\eta-a)]$
and the eigenfunction is within $\arcsin(\rho/\gamma)$ of $\psi$. The
hypothesis is the part that classical presentations leave to the user; making
it a computation is what allows the bound to be used on a mesh-free
approximation whose spectrum is unknown.

The cost is small. In two dimensions a certificate requires four to six
verified counts, each a sparse factorization taking $0.03$\,s at
$81\times81$; in one dimension a count takes $0.4$\,ms, so the count-only
bisection fallback can be used freely and a certificate spends some two
hundred of them. Measured end to end, a full certificate takes $0.19$\,s and
$0.074$\,s against $30$--$110$\,s of training. The guided filter repeats the certificate at each
candidate cutoff---$27$ candidates in the one-dimensional runs and $30$--$62$
in the two-dimensional ones---which costs $2.6$\,s and, scaled from the
measured $16$\,s at $81$ candidates, $6$--$12$\,s, still seconds. There is no reason not to certify every run that is reported.

The benefit is not accuracy but accountability. In our own experiments the
certificate changed decisions four times: it rejected an excited-state run
that had a valid-looking energy (state index mismatch), it exposed a
degenerate state masquerading as a $10^{-12}$-accurate one, it showed that
part of what looked like network error on the finite well was discretization
error that no amount of training could remove, and it decided which pairs of
independently trained networks actually span a degenerate level---accepting
the three pairs that contain the orthogonal partner with bounds of
$5$--$6\times10^{-3}$, refusing the two collapsed pairs, and pricing the
ambiguous pair at $0.68$.

\subsection{Limitations}

\paragraph{The enclosure can be uninformative} For the finite well the ratio
of certified width to true error is $10^{3}$ and the angle bound is $O(1)$.
The bound is correct and useless. This is not a defect of the implementation
but a property of residual-based bounds on non-smooth problems: the residual
cannot distinguish a genuine high-wavenumber component from an error of the
same shape, and the spectral gap at the relevant scale is large enough that
the quadratic Temple gain does not compensate. Any claim that such bounds
``certify PINN accuracy'' must be conditioned on the potential.

\paragraph{Degeneracy is now inside the protocol, but only at counting
resolution} The cluster certificate of Section~\ref{sec:subspace} removes the
structural obstacle: a degenerate level is certified as an invariant subspace,
with the multiplicity measured by the count and the angle bounded by
Theorem~\ref{thm:angle}. Two limitations remain, and they are of different
kinds. The first is the resolution floor \eqref{eq:floor}: two shifts closer
than $10^{-9}s$ cannot be separated by the count, and a cluster whose boundary
margin falls below the floor is refused---including, as the degeneracy
experiment shows, a pair that double precision can barely split. The floor is
a declared safety margin rather than a theorem (Section~\ref{sec:floor}), and
directed rounding does not remove it: it makes the count beneath the margin
exact, so that the refusal happens for a stated reason at a stated pivot,
rather than moving the margin. Counting at a higher precision would move the
floor but not abolish the need for one. The second limitation is structural: a
trial subspace that omits a member of the level is refused because its
dimension disagrees with the count, which is the correct answer but means that
the certificate can detect a wrong basis and cannot repair it. A milder
limitation is that the cluster bound is first order in $\rho$: unlike the
Kato--Temple bound it does not improve quadratically as the residual shrinks,
because a cluster has no interior eigenvalue to isolate. For very accurate
subspaces the counting resolution, not the angle bound, becomes the limiting
quantity.

\paragraph{The certified object is discrete} The gap between $\Hh$ and
$\mathcal H$ is $O(h^2)$ for smooth eigenfunctions and dominates the
certificate at the grid sizes used here (Section~\ref{sec:discretization}). A
certificate for the continuous operator would require either an explicit
$O(h^2)$ bound with computable constants, or a homotopy-based enclosure of the
elliptic problem itself \citep{plum1991}, which needs a priori pointwise
bounds on the eigenfunction that a trained network does not supply; nor is a
$p$-version/spectral discretization with exponential convergence implemented.
We do not claim any of the three. What
Section~\ref{sec:variational} does claim is the one-sided half, and it is
enough to change what a reported energy means: the interpolant of the network
eigenfunction is a conforming trial function, so its Rayleigh quotient bounds
the continuous eigenvalue from above without any regularity input. The other
side is where the difficulty sits, and the finite well shows it is not a
technicality: a certified discrete eigenvalue can lie \emph{above} the
continuous one when the potential is discontinuous, so no care with the
factorization turns a discrete count into a continuous enclosure. The same boundary explains why the count
refers to $\Hh$ and not to the spectrum of a non-conforming trial space: such
a space can carry spurious eigenvalues \citep{lewinsere2010}, so a discrete
eigenvalue computed from it would not by itself identify the state it belongs
to.

\paragraph{Directed rounding in one dimension, not in general two} The count
of Section~\ref{sec:certified} is computed with outward rounding, so in one
dimension---and in two whenever the potential is separable, which covers the
harmonic oscillator and the box---it is a theorem about the stored matrix
rather than a guarded estimate. What remains outside is the general
two-dimensional count for a non-separable potential: the interval elimination
refuses on most shifts at the grid sizes we use, so the anharmonic experiments
still read their count off the floating-point factorization with the checks of
Section~\ref{sec:counting}. We report that measurement rather than route
around it. The separable construction suggests the next step---splitting a
non-separable potential into separable pieces and carrying the remainder as a
perturbation---but we have not implemented it, and we do not claim it here.

\paragraph{Statistical power} We report counts, not rates with confidence
intervals. Every enclosure we report contains the discrete eigenvalue to which
it refers---the eighteen of the main sweep, the nine learnable-energy and nine
reference-energy runs of the leakage audit, and the three deflated runs, for
$39/39$; in the multi-state chain each enclosure contains the eigenvalue of
the state that its certified index names. The Wilson 95\% interval for
$39/39$ is $[0.91,1.0]$, so the data are
still consistent with a per-run failure probability as large as $9\%$; the
eighteen runs of the main sweep alone give $[0.82,1.0]$. The multi-state chain
is the more demanding test and is not perfect: $25/27$ certify the
\emph{requested} state, Wilson $[0.77,0.98]$, with both failures traced to an
upstream state mismatch that the state index flagged. A claim of the form
``the method certifies correctly with probability at least $0.95$'' would
require about seventy runs with no failure, and we do not make it.

\subsection{Why not solve \texorpdfstring{$\Hh$}{H\_h} instead?}
\label{sec:why}

A referee may reasonably ask what the certificate adds, since on every
problem in this paper $\Hh$ is assembled and a shift-invert Lanczos solve of
the same matrix returns $\lambda_h$ and its eigenvector in $20$\,ms at
$81\times81$ and $1.7$\,ms at $N=2001$, roughly ten times faster than the
certificate in two dimensions and forty times in one ($0.19$\,s and
$0.074$\,s; both measured in
\texttt{scripts/bench\_count.py}). We do not think the answer is speed, and
the protocol should not be sold on it. Three things are at stake that a solve of $\Hh$
does not supply.

\emph{The certificate does not need a converged eigenvector.} Counting is the
only operation it performs on $\Hh$, and ARPACK is used to propose intervals
that a count then verifies; the count-only fallback of
Section~\ref{sec:isolating} removes even that. The statement therefore
attaches to the object the user has---one network output, one snapshot, one
subspace---rather than to the object a solver would produce instead.

\emph{At a degenerate level the classical comparison is ill-posed.} The
eigenvector a solver returns for a multiple eigenvalue is an arbitrary element
of the level, chosen by its tie-breaking; the overlap of the same network
output with it can be $1$ or $0$ without anything about the network
changing. The experiments of Section~\ref{sec:deg-results} are exactly this
case: three runs and the solver's eigenvector coincide, the fourth does not,
and every residual-based comparison is silent. The cluster certificate is
invariant under the choice, and it prices the difference. This is the one
place where the certificate states something that \emph{cannot} be read off a
solver's output.

\emph{The audit is of the procedure, not of a number.} State index, two-sided
enclosure, angle bound and measured multiplicity together describe a pipeline
end to end; a comparison with a reference eigenvalue describes one run. The
leakage audit (Section~\ref{sec:leak-results}) and the state-index failures of
the multi-state chain (Section~\ref{sec:multistate}) are decisions of the
first kind. This is also the honest limit of the present work: when $\Hh$
cannot be assembled---the high-dimensional regime that motivates neural
eigen-solvers in the first place---every statement here refers to a matrix
that does not exist, and matrix-free verification is open. The protocol is
addressed to the intermediate regime in which the problem is small enough to
assemble and large enough that the learned function, not the solver, is what
is being reported.

\subsection{Relation to existing work}

Enclosure of eigenvalues by variational bounds plus a verified count is the
standard toolkit of a posteriori eigenvalue analysis for finite elements
\citep{babuska1991}: the Temple and Lehmann--Goerisch bounds are used exactly
as here, and the count is usually obtained from a Cholesky factorization of
the stiffness matrix at a shift. Our contribution is not to that theory but to
its transfer: the approximation is a network, the matrix is assembled for the
certificate alone, and the two failure modes that matter in this setting
(degeneracy and leakage) are given explicit treatments.

For clusters the closest antecedent is \citet{cances2020}, who give
guaranteed two-sided bounds for eigenvalue clusters of conforming
discretizations, multiplicities included. Their analysis is built on the
Galerkin structure---residual estimators, orthogonality, and a comparison
between the approximation space and its neighbourhood. A trained network has
none of these properties, and applying those bounds to one would require
proving an analogue from scratch.
The counting route taken here needs no such hypothesis: it measures the
multiplicity directly and bounds the subspace by the residual-to-margin
ratio, which is why the same statement covers exact, approximate and
accidentally degenerate clusters.

For forward PDEs, a posteriori error estimates for physics-informed networks
have begun to appear \citep{mishra2023}; they bound the error of a solution
operator in a norm, and require regularity assumptions on the data. An
eigenvalue problem differs in that the object to be bounded is a spectral
quantity and the unknown is defined only up to a subspace; the counting
machinery is what replaces the coercivity argument.

Verified-computing libraries such as INTLAB~\citep{rump1999} and
Arb~\citep{johansson2017}, and their Python interfaces, would supply the same
arithmetic at arbitrary precision. We use only IEEE doubles with one-ulp
outward padding (\texttt{nextafter}) because the recurrences need nothing
more: the certified count is a sign test on pivots whose intervals close under
a correctly rounded operation, and the failure mode---a pivot interval that
contains zero---is detected and refused rather than hidden. The same libraries
remain the natural route to a fully rigorous two-dimensional count and to
interval enclosures of the Rayleigh quotients themselves, at a modest
constant-factor cost.

\subsection{A reporting checklist}

The experiments suggest a short checklist for papers that report neural
eigenpairs, each item of which would have changed the interpretation of at
least one run in this study:
\begin{enumerate}
\item state the discrete operator and the grid; report the discretization gap
  separately from the network error;
\item report the certified state index, not only the energy;
\item report the enclosure width, and the ratio of width to true error, not
  only the true error;
\item report raw and filtered quantities side by side, with the $L^2$ change
  caused by the filter;
\item state whether any reference quantity entered the training objective;
\item refuse to certify a degenerate state individually rather than reporting
  its residual as an accuracy, and certify the invariant subspace instead,
  with the multiplicity measured by counting.
\end{enumerate}
Items 2, 3 and 6 are cheap; they require no additional theory and would make
the reported numbers comparable across papers.
